% !TeX root = ../../Book2.tex
\section{特征曲线}
\label{b2:sec:0702}

变系数速度场已经给出粒子在每一点的瞬时运动方向，却没有直接给出整条轨道。先解常微分方程，再用轨道作坐标，是本节的构造顺序。以下 \(U\subseteq\mathbb R^d\) 为开集，\(b\in C^1(U;\mathbb R^d)\)；本节先处理不显含时间的速度场。

\subsection{特征常微分方程}
\label{b2:subsec:070201}

\begin{definition}\label{b2:def:07-characteristic}
设 \(d\) 为正整数，\(U\subset\mathbb R^d\) 为开集，\(b:U\to\mathbb R^d\) 连续，\(I\subset\mathbb R\) 为区间。若 \(X:I\rightarrow U\) 属于 \(C^1\)，并满足
\[
 X'(t)=b\bigl(X(t)\bigr),
\]
则称它为速度场 \(b\) 的\term[tezheng-quxian]{特征曲线}[characteristic curve]，也称特征线。其常微分方程称为\term[tezheng-changweifen-fangcheng]{特征常微分方程}[characteristic ordinary differential equation]。
\end{definition}

沿这样的曲线，复合函数的导数只涉及 \(u_t+b\cdot\nabla u\)。对本节采用的 \(C^1\) 速度场，\cref{b2:thm:07-local-ode}保证经过每个初始点的轨道局部存在且唯一。将其最大解记为 \(\Phi_t(a)\)；流的复合、逆关系和对初始点的 \(C^1\) 依赖已经在\secref{b2:sec:ode-smooth-flow}证明。下面利用这些结论，把沿轨道求出的函数还原为偏微分方程的经典解。

\subsection{沿特征线的表示公式}
\label{b2:subsec:070203}

对自治输运问题 \(u_t+b(x)\cdot\nabla u=0\)，沿轨道有
\[
 \frac{\mathrm d}{\mathrm dt}u\bigl(\Phi_t(a),t\bigr)=0.
\]
因此只要轨道可从时刻 \(t\) 返回 \(0\)，解就由
\begin{equation}\label{b2:eq:07-autonomous-representation}
 u(x,t)=g\bigl(\Phi_{-t}(x)\bigr)
\end{equation}
给出。这里的“返回”既包括时间存在范围，也包括轨道始终位于系数定义域；若在有界区域中轨道先碰到边界，单靠内部初值不足以继续决定数值。

\ProseParagraphBreak

例如平面速度 \(b(x,y)=(-y,x)\) 给出
\[
 \Phi_t(a_1,a_2)=(a_1\cos t-a_2\sin t,\ a_1\sin t+a_2\cos t).
\]
解是初始图形的旋转，\(u(x,y,t)=g(x\cos t+y\sin t,-x\sin t+y\cos t)\)。另一例是 \(b(x)=x\)，其流为 \(\Phi_t(a)=e^ta\)，故 \(u(x,t)=g(e^{-t}x)\)。后者虽然函数值不变，体积却按 \(e^{dt}\) 扩张；把 \(u\) 的积分解释成守恒质量就会出错。

\subsection{全局流与非特征初值问题}
\label{b2:subsec:070204}

\begin{proposition}\label{b2:prop:07-global-flow}
设 \(d\) 为正整数，\(b\in C^1(\mathbb R^d;\mathbb R^d)\)，且存在 \(C\geq0\) 使对所有 \(x\in\mathbb R^d\) 都有
\( |b(x)|\leq C(1+|x|)\)，记 \(\Phi_t\) 为 \(b\) 的自治流。对 \(g\in C^1(\mathbb R^d;\mathbb C)\)，\(u(x,t)=g(\Phi_{-t}(x))\) 给出 \(u_t+b(x)\cdot\nabla u=0\)、\(u(x,0)=g(x)\) 在 \(\mathbb R^d\times\mathbb R\) 上的唯一复值 \(C^1\) 解。
\end{proposition}
\begin{proof}
由\cref{b2:prop:ode-linear-growth}，所有轨道在整个实轴上存在；\cref{b2:thm:07-flow-derivative}给出全局流的 \(C^1\) 性。因此表示公式给出 \(C^1\) 函数。沿特征线该函数等于 \(g(a)\)，故满足方程。任意另一经典解沿同一轨道也恒等于 \(g(a)\)，因此唯一。
\end{proof}

流也给出一般数据曲线的局部求解。把平面上的向量场写成 \(\beta=(\beta_1,\beta_2)\)，给定参数曲线 \(\gamma(s)\) 后，重要的是它是否横穿轨道，而不是曲线能否写成水平直线。

\begin{proposition}[非特征数据曲线]\label{b2:prop:07-noncharacteristic-data}
设 \(U\subset\mathbb R^2\) 为开集，\(\beta\in C^1(U;\mathbb R^2)\)、\(c,f\in C^1(U;\mathbb R)\)，\(I\subset\mathbb R\) 为开区间，\(\gamma\in C^1(I;U)\)、\(g\in C^1(I;\mathbb R)\)，\(s_0\in I\)，且在 \(s=s_0\) 有
\[
 \det\bigl(\gamma'(s_0),\beta(\gamma(s_0))\bigr)\ne0.
\]
则在 \(\gamma(s_0)\) 附近，问题 \(\beta\cdot\nabla u+cu=f\)、\(u(\gamma(s))=g(s)\) 有唯一的局部实值 \(C^1\) 解；这里的曲线数据只对 \(s_0\) 附近的参数规定。局部唯一性指任意两个满足同一曲线数据的局部解，在适当缩小的共同邻域内相等。
\end{proposition}
\begin{proof}
记 \(\Phi_t\) 为 \(\beta\) 的局部流，定义 \(\Psi(s,t)=\Phi_t(\gamma(s))\)。流对时间和初始点为 \(C^1\)，故 \(\Psi\) 也为 \(C^1\)。在 \((s_0,0)\)，它的两列导数正是 \(\gamma'(s_0)\)、\(\beta(\gamma(s_0))\)，故\cref{b2:thm:01-inverse}给出局部 \(C^1\) 逆坐标。缩小到参数矩形 \(J\times(-\tau,\tau)\)，其中 \(J\subset I\) 是含 \(s_0\) 的开区间、\(\tau>0\)，令 \(W=\Psi(J\times(-\tau,\tau))\subset U\)，则 \((s,t)=\Psi^{-1}(x)\) 在 \(W\) 上有定义。沿每个固定 \(s\) 解标量方程
\[
 \partial_tV(s,t)+c(\Psi(s,t))V(s,t)=f(\Psi(s,t)),
 \quad V(s,0)=g(s).
\]
具体地，记
\[
 C(s,t)=\int_0^t c(\Psi(s,r))\,\mathrm dr.
\]
积分因子公式给出唯一的
\[
 V(s,t)=e^{-C(s,t)}\left(g(s)+
       \int_0^t e^{C(s,r)}f(\Psi(s,r))\,\mathrm dr\right).
\]
对参数矩形中的每个紧子矩形，被积函数及一阶导数有共同界，\secref{b2:sec:0103}的积分求导规则于是给出 \(V\in C^1\)。令 \(u=V\circ\Psi^{-1}\)。由于 \(\partial_t\Psi=\beta(\Psi)\)，链式法则给出方程；\(t=0\) 给出曲线数据。任何另一经典解沿同一流也满足这个标量方程，所以在缩小后的坐标邻域内相同。
\end{proof}

设 \(U\subset\mathbb R^2\) 开，\(\beta\in C^1(U;\mathbb R^2)\)，\(\gamma:I\to U\) 为 \(C^1\) 曲线，其中 \(I\) 为开区间。若对每个 \(s\in I\) 都有 \(\det(\gamma'(s),\beta(\gamma(s)))\ne0\)，则称它为\term[fei-tezheng-quxian]{非特征曲线}[noncharacteristic curve]。例如对 \(u_x+u_y=f\) 和 \(\gamma(s)=(s,\kappa s)\)，其中 \(\kappa\in\mathbb R\)，行列式为 \(1-\kappa\)，因此 \(\kappa\ne1\) 时局部数据可独立指定；\(\kappa=1\) 时曲线沿轨道，数据必须满足 \(g'(s)=f(s,s)\)。上述非退化条件使曲线切向和流方向共同组成局部坐标。

\ProseParagraphBreak

\Cref{b2:fig:07-noncharacteristic-data}取同一个常向量场 \(\beta=(1,0)\)，其流为 \(\Phi_t(a_1,a_2)=(a_1+t,a_2)\)。左图取 \(\gamma(s)=(0,s)\)，则 \(\Psi(s,t)=(t,s)\)：数据曲线横穿每条附近的水平轨道，两个参数确实给出二维坐标。右图取 \(\gamma(s)=(s,0)\)，则 \(\Psi(s,t)=(s+t,0)\)：初值参数和流时间都只改变同一条轨道上的位置，不能用它们覆盖二维邻域。

% 本书原创：同一尺度绘出 beta=(1,0) 的精确水平流和两条数据曲线。
\begin{figure}[htbp]
\centering
\begin{tikzpicture}[x=1cm,y=1cm,>=stealth,every node/.style={font=\small},line width=.65pt]
\begin{scope}
\draw[->,gray!60] (-2.2,0)--(2.35,0) node[right] {\(x\)};
\draw[->,gray!60] (0,-1.7)--(0,1.85) node[above] {\(y\)};
\foreach \eta in {-1,0,1}{
  \draw[->,AnalysisGray] (-2.05,\eta)--(2.02,\eta);
}
\draw[->,AnalysisGray,very thick] (.4,1)--(1.4,1) node[midway,above=4pt] {\(\beta=(1,0)\)};
\draw[->,AnalysisBlue,very thick] (0,-1.5)--(0,1.5);
\node[AnalysisBlue,left] at (-.08,1.35) {\(\gamma\)};
\fill[AnalysisBlue] (0,0) circle(1.4pt);
\node at (0,2.6) {横截：\(\gamma(s)=(0,s)\)};
\node at (0,-2.15) {\(\det(\gamma',\beta)=-1\ne0\)};
\end{scope}
\begin{scope}[xshift=6.4cm]
\draw[->,gray!60] (-2.2,0)--(2.35,0) node[right] {\(x\)};
\draw[->,gray!60] (0,-1.7)--(0,1.85) node[above] {\(y\)};
\foreach \eta in {-1,0,1}{
  \draw[->,AnalysisGray] (-2.05,\eta)--(2.02,\eta);
}
\draw[->,AnalysisGray,very thick] (.4,1)--(1.4,1) node[midway,above=4pt] {\(\beta=(1,0)\)};
\draw[->,AnalysisBlue,very thick] (-1.5,0)--(1.5,0);
\node[AnalysisBlue,above] at (-.75,.08) {\(\gamma\)};
\fill[AnalysisBlue] (0,0) circle(1.4pt);
\node at (0,2.6) {相切：\(\gamma(s)=(s,0)\)};
\node at (0,-2.15) {\(\det(\gamma',\beta)=0\)};
\end{scope}
\end{tikzpicture}
\caption[数据曲线与特征轨道的横截和相切]{数据曲线与特征轨道的横截和相切。两图均取 \(\beta=(1,0)\)，灰线为水平轨道，粗灰箭头为长度为 \(1\) 的速度向量，蓝线为数据曲线；两图使用同一坐标尺度。左图的数据曲线与每条附近轨道唯一相交，\(\Psi(s,t)=(t,s)\) 可作局部坐标。右图的数据曲线完全落在一条轨道上，\(\Psi(s,t)=(s+t,0)\) 退化；它仅展示退化可能造成的失败，不代表任意行列式为零的情形都无解或不唯一。}
\label{b2:fig:07-noncharacteristic-data}
\end{figure}


在右图的具体模型中，取方程 \(u_x=0\) 和曲线数据 \(u(s,0)=s\)，沿曲线求导会同时得到 \(u_x(s,0)=1\) 与 \(u_x(s,0)=0\)，所以局部无解。若改取常数数据 \(g\equiv c_0\)，则 \(u(x,y)=c_0+q(y)\) 对任意满足 \(q(0)=0\) 的实值 \(C^1\) 函数 \(q\) 都给出延拓。这个模型说明退化时可能出现的失败方式，并不把一般的行列式为零判作必然无解或必然不唯一；前一命题的局部构造依靠的是逆函数定理的非退化假设。
